\chapter{PHƯƠNG PHÁP ĐỀ XUẤT}

\section{Ký hiệu và các khái niệm cơ bản}

Trước khi đi vào phân tích chi tiết khung lý thuyết NTK, các ký hiệu toán học cơ bản được sử dụng xuyên suốt báo cáo này là:

\begin{enumerate}
    \item \textbf{Vector và ma trận:}
    \begin{itemize}
        \item Vector cột $n$ chiều: $\mathbf{x} \in \mathbb{R}^n$
        \item Phần tử thứ $i$ của vector: $x_i$ hoặc $[\mathbf{x}]_i$
        \item Ma trận đơn vị kích thước $n \times n$: $I_n$
        \item Tích vô hướng: $\langle \mathbf{x}, \mathbf{y} \rangle = \mathbf{x}^\top \mathbf{y} = \sum_{i=1}^{n} x_i y_i$
    \end{itemize}
    
    \item \textbf{Chuẩn Euclidean:} Chuẩn Euclidean (hay chuẩn $L^2$) của một vector $\mathbf{x}$ được định nghĩa là:
    \begin{equation}
        \|\mathbf{x}\| = \sqrt{\sum_{i=1}^{n} x_i^2} = \sqrt{\langle \mathbf{x}, \mathbf{x} \rangle}
    \end{equation}
    
    \item \textbf{Tập hợp chỉ số:} Ký hiệu $[n]$ biểu thị tập hợp các số nguyên dương từ 1 đến $n$:
    \begin{equation}
        [n] = \{1, 2, \ldots, n\}
    \end{equation}
    
    \item \textbf{Bài toán hồi quy và hàm loss:} Xét một mạng nơ-ron với tham số $\theta$ thực hiện ánh xạ đầu vào $\mathbf{x}$ sang đầu ra $F(\mathbf{x}; \theta)$. Cho tập dữ liệu huấn luyện $\{(\mathbf{x}_i, \mathbf{y}_i)\}_{i=1}^{N}$ với $N$ mẫu, mục tiêu huấn luyện là tối thiểu hóa hàm loss Mean Squared Error (MSE):
    \begin{equation}
        \mathcal{L}(\theta) = \frac{1}{2} \|F(\mathbf{x}; \theta) - \mathbf{y}\|^2 = \frac{1}{2} \sum_{i=1}^{N} \|F(\mathbf{x}_i; \theta) - \mathbf{y}_i\|^2
    \end{equation}
    Trong đó $\mathbf{y}$ biểu thị vector chứa tất cả các nhãn đầu ra mục tiêu.
    
    \begin{itemize}
      \item \textbf{Tại sao chọn MSE?}
      \begin{itemize}
          \item Đạo hàm $\frac{\partial \mathcal{L}}{\partial F} = F - \mathbf{y}$ có dạng tuyến tính, cho phép gradient descent hội tụ ổn định.
          \item Trong chế độ NTK, MSE loss dẫn đến nghiệm dạng tường minh theo hồi quy hạt nhân, tạo nền tảng cho phân tích lý thuyết chính xác.
      \end{itemize}
      $\Rightarrow$ Nhờ đó, ta có thể chứng minh rằng dấu của đầu ra dự đoán luôn khớp với nhãn mục tiêu (tức là nếu $y = +1$ thì $\hat{y} > 0$, và ngược lại). Điều này đảm bảo rằng sau khi làm tròn, kết quả sẽ chính xác tuyệt đối - đây là nền tảng cho khái niệm \textit{Khả năng học chính xác} được trình bày trong Định lý 5.1 (Chương 3).
    \end{itemize}
\end{enumerate}

Lệnh \textbf{SBN} có dạng "\textit{trừ và nếu kết quả âm thì nhảy}", nên thỏa cả cập nhật trạng thái, điều kiện, và điều khiển luồng; do đó có thể đạt tính đầy đủ Turing khi được lặp lại/ghép chuỗi phù hợp \cite{minsky_1967,hopcroft_ullman_1979}.

Tác giả đưa SBN vào để nhấn mạnh rằng khung khớp mẫu và phân tích NTK không chỉ nhắm tới các phép toán số học cụ thể, mà còn đóng vai trò \textbf{nền tảng lý thuyết} hướng tới việc học và thực thi \textbf{mọi thuật toán} \cite{de_luca_2025}.

\section{Neural Tangent Kernel (NTK)}

NTK là một công cụ lý thuyết mạnh mẽ để phân tích động lực học huấn luyện của mạng nơ-ron \cite{jacot_2018}. Ý tưởng cốt lõi là mô tả sự thay đổi của đầu ra mạng theo thời gian huấn luyện thông qua một kernel cố định.

\subsection{Ma trận Jacobian}

Xét mạng nơ-ron với đầu ra $F(\mathbf{x}; \theta) \in \mathbb{R}^k$ phụ thuộc vào vector tham số $\theta \in \mathbb{R}^p$. Ma trận Jacobian của đầu ra theo tham số được định nghĩa là:
\begin{equation}
    J(\theta) = \frac{\partial F(\mathbf{x}; \theta)}{\partial \theta} \in \mathbb{R}^{Nk \times p}
\end{equation}

trong đó $N$ là số lượng mẫu dữ liệu, $k$ là chiều của đầu ra, và $p$ là số lượng tham số. Mỗi hàng của ma trận Jacobian biểu thị gradient của một thành phần đầu ra theo toàn bộ tham số.

\subsection{Định nghĩa NTK thực nghiệm}

Neural Tangent Kernel thực nghiệm được định nghĩa thông qua tích của ma trận Jacobian với chuyển vị của nó:
\begin{equation}
    K(\theta) = J(\theta) J(\theta)^\top \in \mathbb{R}^{Nk \times Nk}
\end{equation}

Phần tử $(i,j)$ của ma trận kernel này đo lường mức độ tương quan giữa gradient của đầu ra thứ $i$ và đầu ra thứ $j$ trong không gian tham số. Cụ thể hơn, với hai điểm dữ liệu $\mathbf{x}$ và $\mathbf{x}'$, giá trị kernel được tính như sau:
\begin{equation}
    K(\mathbf{x}, \mathbf{x}'; \theta) = \left\langle \frac{\partial F(\mathbf{x}; \theta)}{\partial \theta}, \frac{\partial F(\mathbf{x}'; \theta)}{\partial \theta} \right\rangle
\end{equation}

\subsection{Ý nghĩa vật lý của NTK}

NTK cho phép hiểu rõ cách mạng nơ-ron \textit{học} từ dữ liệu. Khi huấn luyện bằng gradient descent, sự thay đổi của đầu ra mạng tại một điểm $\mathbf{x}$ không chỉ phụ thuộc vào gradient tại điểm đó mà còn phụ thuộc vào gradient tại tất cả các điểm dữ liệu khác thông qua kernel. Đây chính là cơ chế \textit{học chuyển giao} nội tại trong mạng nơ-ron: thông tin học được từ một mẫu có thể ảnh hưởng đến dự đoán tại các mẫu khác.

\section{Giới hạn chiều rộng vô hạn và Động lực học}

Một trong những đóng góp quan trọng nhất của lý thuyết NTK là khám phá hành vi của mạng nơ-ron khi chiều rộng lớp ẩn tiến tới vô hạn. Trong giới hạn này, việc huấn luyện mạng nơ-ron trở nên đơn giản hóa đáng kể và có thể được phân tích chính xác.

\subsection{Chế độ huấn luyện lười}

Khi chiều rộng lớp ẩn $m \rightarrow \infty$, mạng nơ-ron hoạt động trong chế độ được gọi là \textit{huấn luyện lười}. Trong chế độ này, các tham số của mạng thay đổi cực nhỏ quanh điểm khởi tạo trong suốt quá trình huấn luyện:
\begin{equation}
    \|\theta(t) - \theta(0)\| \approx O\left(\frac{1}{\sqrt{m}}\right) \rightarrow 0 \quad \text{khi } m \rightarrow \infty
\end{equation}

Điều này có vẻ nghịch lý: làm sao mạng có thể học được nếu tham số hầu như không thay đổi? Câu trả lời nằm ở việc với số lượng tham số cực lớn, ngay cả những thay đổi nhỏ trong mỗi tham số cũng đủ để tạo ra thay đổi đáng kể trong đầu ra. Đây chính là sức mạnh của việc có nhiều chiều trong không gian tham số.

\subsection{Sự ổn định của Kernel (Kernel đóng băng)}

Một hệ quả quan trọng của huấn luyện lười là NTK trở nên ổn định (hay đóng băng) trong suốt quá trình huấn luyện. Cụ thể, khi $m \rightarrow \infty$:
\begin{equation}
    K_m(\theta(t)) \rightarrow \Theta(\mathbf{x}, \mathbf{x}') \quad \text{với mọi } t \geq 0
\end{equation}

trong đó $\Theta(\mathbf{x}, \mathbf{x}')$ là \textit{NTK với chiều rộng lớp ẩn vô hạn}, một kernel xác định và cố định chỉ phụ thuộc vào cặp đầu vào $(\mathbf{x}, \mathbf{x}')$ và kiến trúc mạng, không phụ thuộc vào quá trình huấn luyện. Sự ổn định này là chìa khóa cho phép phân tích chính xác động lực học huấn luyện.

\subsection{Phương trình vi phân cho đầu ra mạng}

Xét việc huấn luyện mạng bằng gradient flow (gradient descent với learning rate vô cùng nhỏ). Động lực học của tham số được mô tả bởi:
\begin{equation}
    \frac{d\theta}{dt} = -\nabla_\theta \mathcal{L}(\theta) = -J(\theta)^\top (F(\mathbf{x}; \theta) - \mathbf{y})
\end{equation}

Từ đây, có thể suy ra phương trình vi phân cho đầu ra mạng. Đặt $f(t) = F(\mathbf{x}; \theta(t))$ là vector đầu ra tại thời điểm $t$, ta có:
\begin{equation}
    \frac{df}{dt} = J(\theta) \frac{d\theta}{dt} = -J(\theta) J(\theta)^\top (f(t) - \mathbf{y}) = -K(\theta(t))(f(t) - \mathbf{y})
\end{equation}

Trong giới hạn chiều rộng vô hạn, với $K(\theta(t)) \approx \Theta$ cố định, phương trình trên trở thành phương trình vi phân tuyến tính:
\begin{equation}
    \frac{df}{dt} = -\Theta(f(t) - \mathbf{y})
\end{equation}

Nghiệm của phương trình này có dạng giải tích tường minh:
\begin{equation}
    f(t) = \mathbf{y} + e^{-\Theta t}(f(0) - \mathbf{y})
\end{equation}

trong đó $e^{-\Theta t}$ là hàm mũ của ma trận. Khi $t \rightarrow \infty$, nếu $\Theta$ là ma trận xác định dương, ta có $f(t) \rightarrow \mathbf{y}$, nghĩa là mạng học được chính xác nhãn đầu ra.

\subsection{Tương đương với Hồi quy hạt nhân}

Kết quả quan trọng nhất của phân tích trên là việc huấn luyện mạng nơ-ron trong giới hạn chiều rộng vô hạn hoàn toàn tương đương với bài toán \textit{Hồi quy hạt nhân} cổ điển. Cho một điểm test mới $\mathbf{x}^*$, dự đoán của mạng sau khi huấn luyện hội tụ là:
\begin{equation}
    f_\infty(\mathbf{x}^*) = \Theta(\mathbf{x}^*, \mathcal{X}) \Theta(\mathcal{X}, \mathcal{X})^{-1} \mathbf{y}
\end{equation}

trong đó:
\begin{quote}
\begin{itemize}
    \item $\mathcal{X} = \{\mathbf{x}_1, \ldots, \mathbf{x}_N\}$ là tập các điểm huấn luyện
    \item $\Theta(\mathbf{x}^*, \mathcal{X}) \in \mathbb{R}^{1 \times N}$ là vector kernel giữa điểm test và các điểm huấn luyện
    \item $\Theta(\mathcal{X}, \mathcal{X}) \in \mathbb{R}^{N \times N}$ là ma trận kernel (ma trận Gram) trên tập huấn luyện
\end{itemize}
\end{quote}

$\Rightarrow$ Công thức này chính là nghiệm của bài toán Hồi quy hạt nhân tiêu chuẩn. Điều này cho thấy mạng nơ-ron sâu, trong giới hạn chiều rộng vô hạn, hoạt động như một phương pháp kernel cổ điển với một kernel đặc biệt được xác định bởi kiến trúc mạng.

\section{Kiến trúc mạng nơ-ron hai lớp}

Nghiên cứu tập trung vào mạng nơ-ron hai lớp với kiến trúc đơn giản nhưng đủ mạnh để phân tích lý thuyết.

\subsection{Cấu trúc mạng}

Mạng nơ-ron hai lớp được sử dụng có cấu trúc như sau:
\begin{quote}
\begin{itemize}
    \item \textbf{Lớp đầu vào:} Nhận vector $\mathbf{x} \in \mathbb{R}^d$
    \item \textbf{Lớp ẩn:} Gồm $m$ nơ-ron với hàm kích hoạt ReLU, không sử dụng hệ số chệch
    \item \textbf{Lớp đầu ra:} Ánh xạ tuyến tính từ $m$ nơ-ron ẩn sang $k$ đầu ra
\end{itemize}
\end{quote}

Công thức toán học của mạng được viết là:
\begin{equation}
    F(\mathbf{x}) = W^2 \cdot \text{ReLU}(W^1 \mathbf{x})
\end{equation}

trong đó:
\begin{quote}
\begin{itemize}
    \item $W^1 \in \mathbb{R}^{m \times d}$ là ma trận trọng số lớp thứ nhất
    \item $W^2 \in \mathbb{R}^{k \times m}$ là ma trận trọng số lớp thứ hai
    \item $\text{ReLU}(z) = \max(0, z)$ là hàm kích hoạt ReLU được áp dụng trên từng phần tử
\end{itemize}
\end{quote}

Việc không sử dụng hệ số chệch đơn giản hóa phân tích lý thuyết mà không làm mất tính tổng quát về mặt biểu diễn, vì hệ số chệch có thể được mô phỏng bằng cách thêm một chiều hằng số vào đầu vào.

\subsection{Tham số hóa NTK}

Để đảm bảo tính hội tụ trong giới hạn chiều rộng vô hạn, các trọng số được khởi tạo theo chuẩn \textit{Tham số hóa NTK}. Cụ thể:
\begin{quote}
\begin{itemize}
    \item Các phần tử của $W^1$ được khởi tạo i.i.d. từ phân phối $\mathcal{N}(0, 1)$
    \item Các phần tử của $W^2$ được khởi tạo i.i.d. từ phân phối $\mathcal{N}(0, 1/m)$
\end{itemize}
\end{quote}

Hệ số $1/m$ trong phương sai của $W^2$ là then chốt. Nó đảm bảo rằng đầu ra của mạng có phương sai $O(1)$ khi $m \rightarrow \infty$:
\begin{equation}
    \text{Var}[F(\mathbf{x})] = \text{Var}\left[\sum_{j=1}^{m} W^2_{ij} \cdot \text{ReLU}(W^1_j \mathbf{x})\right] \approx m \cdot \frac{1}{m} \cdot O(1) = O(1)
\end{equation}

Nếu không có hệ số này, đầu ra sẽ có phương sai tăng theo $m$, làm mất tính ổn định khi $m \rightarrow \infty$.

\section{Điều kiện về khả năng học chính xác}

\subsection{Giả thuyết 5.1: Kiểm soát tương quan không mong muốn}

Xét một bước tính toán của thuật toán với đầu vào được biểu diễn dưới dạng mẫu. Đầu ra của bộ dự đoán NTK có thể được phân tích thành hai thành phần:
\begin{equation}
    \hat{y} = w^1 \cdot (\text{tín hiệu chính xác}) + w^0 \cdot (\text{tổng các tương quan không mong muốn})
\end{equation}

trong đó:
\begin{quote}
\begin{itemize}
    \item $w^1$ là \textit{trọng số tín hiệu}: Đóng góp từ các mẫu khớp chính xác với quy tắc cần áp dụng
    \item $w^0$ là \textit{trọng số nhiễu}: Đóng góp từ các mẫu không liên quan nhưng có tương quan khác không với đầu vào
\end{itemize}
\end{quote}

	\textbf{Định nghĩa tương quan không mong muốn:} Một tương quan được gọi là không mong muốn khi một mẫu trong tập huấn luyện có tích vô hướng khác không với đầu vào hiện tại, mặc dù mẫu đó không đại diện cho quy tắc cần áp dụng. Gọi $C$ là số lượng các tương quan không mong muốn như vậy.

\textbf{Điều kiện đủ để học chính xác:} Để đầu ra $\hat{y}$ có dấu đúng (từ đó có thể làm tròn về giá trị chính xác), cần đảm bảo rằng thành phần tín hiệu lấn át thành phần nhiễu:
\begin{equation}
    |w^1| > C \cdot |w^0|
\end{equation}

hay tương đương:
\begin{equation}
    C < \frac{|w^1|}{|w^0|}
\end{equation}

Điều kiện này có ý nghĩa trực quan: số lượng tương quan không mong muốn phải đủ nhỏ so với tỷ lệ giữa trọng số tín hiệu và trọng số nhiễu. Trong thiết kế mẫu của nghiên cứu này, các tác giả đã chứng minh rằng với cách mã hóa phù hợp, tỷ lệ $|w^1|/|w^0|$ tăng theo kích thước bài toán, trong khi $C$ chỉ tăng theo logarithm, đảm bảo điều kiện trên được thỏa mãn \cite{de_luca_2025}.

\subsection{Vai trò của hàm làm tròn}

Một thành phần quan trọng trong khung lý thuyết là việc áp dụng hàm làm tròn sau mỗi bước dự đoán. Cụ thể:
\begin{equation}
    y_{\text{output}} = \text{round}(\hat{y}) = 
    \begin{cases}
        +1 & \text{nếu } \hat{y} > 0 \\
        -1 & \text{nếu } \hat{y} < 0
    \end{cases}
\end{equation}

Hàm làm tròn đóng vai trò then chốt trong việc:
\begin{enumerate}
    \item \textbf{Loại bỏ nhiễu:} Miễn là đầu ra thô $\hat{y}$ có dấu đúng, giá trị chính xác được khôi phục hoàn toàn sau khi làm tròn. Điều này cho phép chấp nhận một lượng nhiễu nhất định mà không ảnh hưởng đến kết quả cuối cùng.
    
    \item \textbf{Ngăn chặn tích lũy sai số:} Trong các thuật toán đa bước, nếu không có làm tròn, sai số nhỏ ở các bước trước sẽ tích lũy và phóng đại qua các bước sau. Làm tròn tại mỗi bước đảm bảo rằng đầu ra luôn chính xác tuyệt đối, ngăn chặn hiện tượng này.
    
    \item \textbf{Chuyển đổi xấp xỉ thành chính xác:} Trong khi bộ dự đoán NTK về bản chất là một bộ xấp xỉ liên tục, việc kết hợp với làm tròn biến nó thành một hệ thống thực thi chính xác các phép toán rời rạc.
\end{enumerate}

Kết hợp Giả thuyết 5.1 và cơ chế làm tròn, nghiên cứu chứng minh rằng mạng nơ-ron trong giới hạn chiều rộng vô hạn có khả năng học và thực thi chính xác các chỉ thị thuật toán, mở đường cho các kết quả về phép cộng, phép nhân, và chỉ thị SBN được trình bày trong các chương tiếp theo.

\section{Kỹ thuật Khớp mẫu}

Ý tưởng cốt lõi của phương pháp đề xuất là phân rã thuật toán thành các mẫu cục bộ, mỗi mẫu thao tác trên một khối bit nhỏ. Thay vì yêu cầu mạng nơ-ron học toàn bộ hàm phức tạp, ta chuyển bài toán thành việc nhận dạng và áp dụng các quy tắc cục bộ đơn giản.

\subsection{Nguyên lý phân rã thuật toán}

Xét một thuật toán hoạt động trên chuỗi bit $\mathbf{x} = (x_1, x_2, \ldots, x_l)$ với $l$ bit. Thuật toán được phân rã thành $n$ khối, mỗi khối $i$ chịu trách nhiệm tính toán một phần đầu ra dựa trên một tập con các bit đầu vào. Mỗi khối được đặc trưng bởi một hàm cục bộ $f_i$ thao tác trên các bit trong phạm vi của nó.

\vspace{0.2cm}

\textbf{Hàm khớp mẫu cục bộ:} Mỗi hàm cục bộ $f_i: \{-1, +1\}^{b_i} \rightarrow \{-1, +1\}^{o_i}$ nhận đầu vào là $b_i$ bit và trả về $o_i$ bit đầu ra. Hàm này được định nghĩa thông qua một tập các mẫu:
\begin{equation}
    f_i(\mathbf{x}_{[S_i]}) = \bigvee_{j \in T_i} \mathbf{1}[\mathbf{x}_{[S_i]} = \mathbf{t}_j] \cdot \mathbf{o}_j
\end{equation}

trong đó:
\begin{quote}
\begin{itemize}
    \item $S_i$ là tập chỉ số các bit đầu vào của khối $i$.
    \item $\mathbf{x}_{[S_i]}$ là vector con của $\mathbf{x}$ tại các vị trí trong $S_i$.
    \item $T_i$ là tập các mẫu thuộc khối $i$.
    \item $\mathbf{t}_j$ là vector mẫu đầu vào thứ $j$.
    \item $\mathbf{o}_j$ là vector đầu ra tương ứng với mẫu $j$.
    \item $\mathbf{1}[\cdot]$ là hàm chỉ thị.
\end{itemize}
\end{quote}

\vspace{0.2cm}

\textbf{Hàm tổng hợp toàn cục:} Đầu ra của thuật toán được tổng hợp từ các hàm cục bộ thông qua phép tuyển theo bit (OR theo bit):
\begin{equation}
    f(\hat{\mathbf{x}}) = \bigvee_{i=1}^{n} f_i(\hat{\mathbf{x}}_{[S_i]})
\end{equation}

Phép tuyển $\bigvee$ hoạt động như phép OR trên các bit: nếu bất kỳ hàm cục bộ nào trả về bit 1 tại một vị trí, bit đầu ra tại vị trí đó sẽ là 1.

\subsection{Cơ chế thực thi lặp}

Nhiều thuật toán yêu cầu thực thi qua nhiều bước lặp, trong đó trạng thái của bước tiếp theo phụ thuộc vào kết quả của bước hiện tại. Phương pháp khớp mẫu hỗ trợ điều này thông qua cơ chế thực thi lặp:
\begin{equation}
    \hat{\mathbf{x}}_{t+1} = f(\hat{\mathbf{x}}_t)
\end{equation}

trong đó $\hat{\mathbf{x}}_t$ là trạng thái tại bước thứ $t$. Quá trình lặp tiếp tục cho đến khi đạt được trạng thái kết thúc hoặc sau một số bước cố định được xác định bởi thuật toán.

Điểm quan trọng là sau mỗi bước lặp, đầu ra được làm tròn về các giá trị rời rạc $\{-1, +1\}$ trước khi đưa vào bước tiếp theo. Điều này ngăn chặn tích lũy sai số và đảm bảo tính chính xác tuyệt đối.

\section{Cấu trúc dữ liệu và Trực giao hóa}

Để huấn luyện mạng nơ-ron học được các mẫu, cần có chiến lược mã hóa phù hợp biến đổi các quy tắc thuật toán thành dữ liệu huấn luyện.

\subsection{Mô hình Bag of Rule}

Chiến lược mã hóa được sử dụng là gom tất cả các quy tắc cục bộ của thuật toán vào một không gian vector chung, gọi là mô hình Bag of Rules. Mỗi quy tắc cục bộ được biểu diễn như một cặp (đầu vào, đầu ra) trong không gian này.

Cụ thể, với thuật toán có $n$ khối và tổng cộng $M$ quy tắc cục bộ, tập huấn luyện có dạng:
\begin{equation}
    \mathcal{D} = \{(\mathbf{q}_j, y_j)\}_{j=1}^{M}
\end{equation}

trong đó $\mathbf{q}_j$ là vector mã hóa quy tắc thứ $j$ và $y_j$ là nhãn đầu ra tương ứng.

\subsection{Mã hóa phân vùng khối}

Mỗi mẫu được mã hóa thành vector $\mathbf{q}_{ij}$ theo cấu trúc phân vùng khối. Vector này có dạng:
\begin{equation}
    \mathbf{q}_{ij} = (\underbrace{0, \ldots, 0}_{\text{khối } 1}, \ldots, \underbrace{\mathbf{t}_{ij}}_{\text{khối } i}, \ldots, \underbrace{0, \ldots, 0}_{\text{khối } n})
\end{equation}

Ý nghĩa của cấu trúc này:
\begin{quote}
\begin{itemize}
    \item Vector được chia thành $n$ phân đoạn, mỗi phân đoạn tương ứng với một khối của thuật toán.
    \item Chỉ phân đoạn thứ $i$ chứa thông tin mẫu $\mathbf{t}_{ij}$, các phân đoạn khác đều bằng 0.
    \item Điều này đảm bảo các mẫu thuộc các khối khác nhau có tích vô hướng bằng 0 với nhau.
\end{itemize}
\end{quote}

\subsection{Trực giao hóa và Ma trận NTK}

Tầm quan trọng của cấu trúc phân vùng khối nằm ở việc nó tạo ra tính trực giao trong dữ liệu huấn luyện. Khi các mẫu thuộc các khối khác nhau trực giao với nhau, ma trận NTK có cấu trúc đặc biệt.

Gọi $\Theta$ là ma trận NTK trên tập huấn luyện. Với dữ liệu được mã hóa phân vùng khối, ma trận này tiệm cận dạng:
\begin{equation}
    \Theta \approx \alpha I_M + \beta \mathbf{1}\mathbf{1}^\top + \text{(nhiễu bậc thấp)}
\end{equation}

trong đó:
\begin{quote}
\begin{itemize}
    \item $\alpha$ là hệ số đường chéo chính.
    \item $I_M$ là ma trận đơn vị kích thước $M \times M$.
    \item $\beta \mathbf{1}\mathbf{1}^\top$ là ma trận nhiễu hạng 1.
    \item $\mathbf{1}$ là vector toàn 1.
\end{itemize}
\end{quote}

Cấu trúc \textbf{ma trận đơn vị + nhiễu hạng 1} này cho phép phân tích chính xác hành vi của bộ dự đoán NTK và là nền tảng cho các định lý về khả năng học chính xác.

\section{Định lý về bộ dự báo NTK}

\subsection{Định lý 5.1: Bảo toàn thông tin dấu}

    	\textbf{Định lý 5.1 (Hành vi của bộ dự báo NTK):} Cho mạng nơ-ron hai lớp trong chế độ NTK ở giới hạn chiều rộng vô hạn được huấn luyện trên tập dữ liệu mẫu $\mathcal{D}$. Gọi $\mu(\hat{\mathbf{x}})$ là kỳ vọng của phân phối giới hạn của đầu ra mạng tại điểm kiểm thử $\hat{\mathbf{x}}$. Khi đó, với mỗi bit đầu ra thứ $i$:
\begin{equation}
    \text{sign}(\mu(\hat{\mathbf{x}})_i) = y_i^*
\end{equation}

trong đó $y_i^* \in \{-1, +1\}$ là giá trị bit đầu ra đúng theo thuật toán.

\subsection{Ý nghĩa của Định lý}

Định lý 5.1 khẳng định rằng bộ dự đoán NTK bảo toàn thông tin về dấu của bit đầu ra:
\begin{itemize}
    \item Nếu bit đầu ra đúng là $+1$, thì $\mu(\hat{\mathbf{x}})_i > 0$.
    \item Nếu bit đầu ra đúng là $-1$, thì $\mu(\hat{\mathbf{x}})_i < 0$.
\end{itemize}

Điều này có ý nghĩa thực tiễn quan trọng: mặc dù đầu ra thô của mạng là một giá trị thực (có thể không chính xác bằng $\pm 1$), việc áp dụng hàm làm tròn sẽ khôi phục giá trị chính xác:
\begin{equation}
    y_{\text{output}} = \text{round}(\mu(\hat{\mathbf{x}})_i) = \text{sign}(\mu(\hat{\mathbf{x}})_i) = y_i^*
\end{equation}

Kết quả này cho phép mạng nơ-ron đạt được độ chính xác tuyệt đối thay vì chỉ xấp xỉ, giải quyết thách thức cốt lõi về việc học các hàm rời rạc.

\section{Kiểm soát tương quan không mong muốn}

Định lý 5.1 không tự động đúng cho mọi thuật toán. Để đảm bảo tính đúng đắn, cần kiểm soát các tương quan không mong muốn trong dữ liệu huấn luyện.

\subsection{Giả thuyết 5.1: Điều kiện về tương quan}

Phân tích đầu ra của bộ dự đoán NTK cho thấy:
\begin{equation}
    \mu(\hat{\mathbf{x}})_i = w^1 \cdot y_i^* + w^0 \cdot \sum_{j \in U_i} y_j
\end{equation}

trong đó:
\begin{quote}
\begin{itemize}
    \item $w^1 > 0$ là \textit{trọng số tín hiệu}: Đóng góp từ mẫu khớp chính xác.
    \item $w^0$ là \textit{trọng số nhiễu}: Đóng góp từ các mẫu có tương quan khác không.
    \item $U_i$ là tập các chỉ số của các mẫu không khớp nhưng có tích vô hướng khác không với đầu vào.
    \item $y_j$ là nhãn của mẫu thứ $j$ trong tập nhiễu.
\end{itemize}
\end{quote}

        \textbf{Giả thuyết 5.1:} Gọi $C_i = |U_i|$ là số lượng tương quan không mong muốn tại bit thứ $i$. Điều kiện để đầu ra có dấu đúng là:
\begin{equation}
    C_i < -\frac{w^1}{w^0}
\end{equation}

Lưu ý rằng $w^0 < 0$ trong các trường hợp được xét, do đó $-w^1/w^0 > 0$.

\subsection{Phân tích trường hợp xấu nhất}

Trong trường hợp xấu nhất, tất cả các tương quan không mong muốn đều đồng pha (có cùng dấu với nhiễu), tổng nhiễu có thể đạt:
\begin{equation}
    \left|\sum_{j \in U_i} y_j\right| \leq C_i
\end{equation}

Để đảm bảo thành phần tín hiệu lấn át nhiễu:
\begin{equation}
    |w^1 \cdot y_i^*| > |w^0 \cdot C_i| \Rightarrow |w^1| > C_i \cdot |w^0| \Rightarrow C_i < \frac{|w^1|}{|w^0|}
\end{equation}

\subsection{Kết quả cho các thuật toán cụ thể}

Các tác giả đã chứng minh rằng các thuật toán được xét đều thỏa mãn giả thuyết 5.1:
\begin{itemize}
    \item \textbf{Phép hoán vị:} $C_i = 0$ (không có tương quan không mong muốn).
    \item \textbf{Phép cộng:} $C_i \leq 1$ (tối đa 1 tương quan không mong muốn).
    \item \textbf{Phép nhân:} $C_i = O(\log l)$ với $l$ là số bit.
    \item \textbf{SBN:} $C_i = O(\log l)$.
\end{itemize}

Với thiết kế mẫu phù hợp, tỷ lệ $|w^1|/|w^0|$ tăng nhanh hơn $C_i$ khi kích thước bài toán tăng, đảm bảo điều kiện được thỏa mãn.

\section{Triển khai thuật toán cụ thể}

\subsection{Phép hoán vị}

Phép hoán vị $\pi: [l] \rightarrow [l]$ ánh xạ vị trí bit $i$ sang vị trí mới $\pi(i)$. Đây là thuật toán đơn giản nhất để minh họa phương pháp.

\vspace{0.2cm}

\textbf{Mã hóa mẫu:} Mỗi bit được xử lý độc lập. Với bit thứ $i$, có 2 mẫu:
\begin{itemize}
    \item Mẫu 1: Đầu vào $x_i = +1 \Rightarrow$ Đầu ra tại vị trí $\pi(i)$ là $+1$.
    \item Mẫu 2: Đầu vào $x_i = -1 \Rightarrow$ Đầu ra tại vị trí $\pi(i)$ là $-1$.
\end{itemize}

\textbf{Số lượng mẫu:} $2l$ mẫu tổng cộng.

\vspace{0.2cm}

\textbf{Tương quan không mong muốn:} Vì mỗi mẫu chỉ phụ thuộc vào một bit đầu vào duy nhất và các bit được mã hóa trong các khối riêng biệt, không có tương quan không mong muốn ($C_i = 0$).

\subsection{Phép cộng}

Phép cộng hai số nhị phân $l$-bit được triển khai thông qua mô phỏng bộ cộng ripple-carry. Thuật toán gồm 2 giai đoạn:

\vspace{0.2cm}

\textbf{Giai đoạn 1 - Cộng theo bit:} Tính tổng từng bit và bit nhớ cục bộ:
\begin{align}
    s_i &= a_i \oplus b_i \oplus c_{i-1} \\
    c_i &= \text{MAJ}(a_i, b_i, c_{i-1})
\end{align}

trong đó $\oplus$ là phép XOR, MAJ là hàm majority, và $c_{i-1}$ là bit nhớ từ vị trí trước.

\vspace{0.2cm}

\textbf{Giai đoạn 2 - Lan truyền bit nhớ:} Bit nhớ được lan truyền từ phải sang trái. Quá trình này yêu cầu $O(\log l)$ bước lặp do kỹ thuật tiền tố song song.

\vspace{0.2cm}

\textbf{Mẫu cho mỗi vị trí:} Mỗi vị trí bit $i$ có $2^3 = 8$ mẫu (tương ứng với 8 tổ hợp của $a_i, b_i, c_{i-1}$).

\vspace{0.2cm}

\textbf{Tương quan không mong muốn:} Tối đa $C_i = 1$, xảy ra khi hai mẫu có cùng giá trị $(a_i, b_i)$ nhưng khác $c_{i-1}$.

\subsection{Phép nhân}

Phép nhân hai số nhị phân $l$-bit sử dụng thuật toán dịch và cộng cổ điển. Thuật toán lặp qua từng bit của số nhân và thực hiện:
\begin{enumerate}
    \item Kiểm tra bit phải cùng của số nhân.
    \item Nếu bit phải cùng = 1: Cộng số bị nhân vào kết quả tích lũy.
    \item Dịch trái số bị nhân và dịch phải số nhân.
\end{enumerate}

\textbf{Số bước lặp:} $l$ bước cho phép nhân $l$-bit.

\vspace{0.3cm}

\textbf{Mã hóa mẫu:} Mỗi bước lặp được mã hóa thành các mẫu riêng, bao gồm:

\begin{itemize}
    \item Mẫu kiểm tra bit phải cùng.
    \item Mẫu cộng có điều kiện.
    \item Mẫu dịch bit.
\end{itemize}

\vspace{-0.1cm}

\textbf{Tương quan không mong muốn:} $C_i = O(\log l)$, phát sinh từ sự chồng chéo giữa các bước lặp.

\subsection{Chỉ thị SBN (Subtract and Branch if Negative)}

SBN là một chỉ thị máy tính đơn giản nhưng đủ mạnh để đạt tính đầy đủ Turing. Chỉ thị SBN(A, B, C) thực hiện:
\begin{enumerate}
    \item Tính $\text{Mem}[A] \leftarrow \text{Mem}[A] - \text{Mem}[B]$
    \item Nếu kết quả âm, nhảy đến địa chỉ C; ngược lại, tiếp tục tuần tự.
\end{enumerate}

\textbf{Tính đầy đủ Turing:} SBN được chứng minh là đầy đủ Turing, nghĩa là bất kỳ hàm tính toán được nào cũng có thể được biểu diễn bằng một chuỗi hữu hạn các chỉ thị SBN.

Mã hóa mẫu cho SBN: Bao gồm các thành phần:
\begin{itemize}
    \item Mẫu cho phép trừ: Tương tự phép cộng với số bù 2.
    \item Mẫu kiểm tra dấu kết quả.
    \item Mẫu cập nhật bộ đếm chương trình.
    \item Mẫu cho phép trừ: Tương tự phép cộng với số bù 2.
    \item Mẫu kiểm tra dấu kết quả.
    \item Mẫu cập nhật bộ đếm chương trình.
\end{itemize}

\textbf{Ý nghĩa:} Việc chứng minh mạng nơ-ron có thể học chính xác SBN đồng nghĩa với việc mạng nơ-ron có khả năng, về nguyên tắc, học thực thi bất kỳ thuật toán nào có thể được mô tả bằng chương trình máy tính.

\section{Độ phức tạp tập hợp}

Mặc dù lý thuyết NTK đảm bảo tính đúng đắn về kỳ vọng, phương sai của từng mô hình riêng lẻ có thể ảnh hưởng đến kết quả thực tế. Để giảm thiểu rủi ro này, phương pháp sử dụng kỹ thuật kết hợp nhiều mô hình.

\subsection{Tại sao cần kỹ thuật kết hợp nhiều mô hình?}

Trong thực tế, với mạng có chiều rộng hữu hạn $m$, đầu ra có phân phối:
\begin{equation}
    F(\hat{\mathbf{x}}) \sim \mathcal{N}(\mu(\hat{\mathbf{x}}), \sigma^2(\hat{\mathbf{x}}))
\end{equation}

trong đó $\sigma^2$ là phương sai. Ngay cả khi $\mu$ có dấu đúng, phương sai lớn có thể khiến một số lần thực thi cho kết quả sai.

\vspace{0.2cm}

Kỹ thuật giải quyết vấn đề này bằng cách huấn luyện $N$ mô hình độc lập và lấy trung bình:
\begin{equation}
    \bar{F}(\hat{\mathbf{x}}) = \frac{1}{N} \sum_{k=1}^{N} F_k(\hat{\mathbf{x}})
\end{equation}

Phương sai của đầu ra trung bình giảm theo $1/N$:
\begin{equation}
    \text{Var}[\bar{F}] = \frac{\sigma^2}{N}
\end{equation}

\subsection{Công thức tính số lượng mô hình tối thiểu}

Để đảm bảo xác suất lỗi không vượt quá $\delta$, số lượng mô hình $N$ cần thỏa mãn (Công thức 8 trong paper):
\begin{equation}
    N \geq \frac{2\sigma^2}{\mu^2} \ln\left(\frac{2}{\delta}\right)
\end{equation}

trong đó:
\begin{quote}
\begin{itemize}
    \item $\mu$ là kỳ vọng của đầu ra (đã được chứng minh có dấu đúng).
    \item $\sigma^2$ là phương sai của từng mô hình.
    \item $\delta$ là xác suất lỗi cho phép.
\end{itemize}
\end{quote}

Công thức này dựa trên các bất đẳng thức tập trung cho tổng các biến ngẫu nhiên Gaussian độc lập.

\subsection{Độ phức tạp theo kích thước bài toán}

Phân tích cho thấy tỷ lệ $\sigma^2/\mu^2$ tăng theo $O(l^2)$ với $l$ là số bit của bài toán. Do đó, độ phức tạp ensemble là:
\begin{equation}
    N = O(l^2 \log(1/\delta))
\end{equation}

Kết quả này có ý nghĩa thực tiễn quan trọng:
\begin{itemize}
    \item Độ phức tạp chỉ tăng đa thức theo kích thước bài toán.
    \item Với $l = 64$ bit (số nguyên 64-bit), cần khoảng $N = O(4096)$ mô hình.
    \item Các mô hình có thể được huấn luyện song song, phù hợp với phần cứng GPU hiện đại.
\end{itemize}

Kết luận, phương pháp đề xuất có khả năng mở rộng tốt cho các bài toán thực tế, không bị giới hạn bởi sự tăng trưởng hàm mũ thường gặp trong các phương pháp học máy truyền thống cho bài toán tổ hợp.
